% Part: lambda-calculus
% Chapter: lambda-definability
% Section: pairs

\documentclass[../../../include/open-logic-section]{subfiles}

\begin{document}

\olfileid[pt]{lam}{ldf}{pai}
\olsection{Pares e predecessor}

\begin{defn}
O par de $M$ e $N$ (escrito $\tuple{M,N}$) é definido da seguinte maneira:
\[
\tuple{M,N} \ident \lambd[f][fMN].
\]
\end{defn}
  
Intuitivamente, trata-se de uma função que recebe uma função e aplica essa
função aos dois elementos do par. Seguindo essa ideia, temos este
construtor, que recebe dois termos e devolve o par que os contém:
\[
  \fn{Pair} \ident \lambd[mn][\lambd[f][fmn]]
\]
Dado um par, também queremos recuperar seus elementos.
Para isso, precisamos de duas funções de acesso, que recebem um par como argumento e
devolvem o primeiro ou o segundo elemento dele:
\begin{align*}
  \fn{Fst} & \ident \lambd[p][p(\lambd[mn][m])]\\
  \fn{Snd} & \ident \lambd[p][p(\lambd[mn][n])]
\end{align*}

\begin{prob}
  Explique por que as funções de acesso $\fn{Fst}$ e $\fn{Snd}$ funcionam.
\end{prob}

Agora, com pares, podemos !!{lambda define} a função predecessor:
\[
  \fn{Pred} \ident \lambd[n][\fn{Fst}(n (\lambd[p][\tuple{\fn{Snd}\, {p}, \fn{Succ}(\fn{Snd}\, {p})}]) \tuple{\num 0, \num 0})]
\]
Lembre-se de que $\num n\, f x$ se reduz a $f^{n}(x)$; neste
caso, $f$ é uma função que recebe um par $p$ e devolve um novo
par contendo o segundo componente de $p$ e o sucessor do
segundo componente; $x$ é o par $\tuple{0,0}$. Assim, o
resultado é $\tuple{0,0}$ para $n=0$, e $\tuple{\num{n-1}, \num n}$
nos demais casos. $\fn{Pred}$ devolve então o primeiro componente do resultado.

A subtração pode ser definida como $\fn{Pred}$ aplicada a $a$, $b$ vezes:
\[
\fn{Sub} \ident \lambd[ab][b \fn{Pred}\, a].
\]
    
\end{document}
